\section{Устройство робота}

Робот обкачивает два источника корпуса \enquote{кино и сериалы}: русскую Википедию
(статьи о фильмах, сериалах и людях кино) и film.ru (рецензии, подборки, интервью).
Работа состоит из двух программ. Обе принимают единственный аргумент --- путь до
общего yaml-конфига.

\begin{itemize}
    \item \texttt{src/crawler.py} --- обкачка: берёт URL из очереди, скачивает
          документ, выделяет текст, сохраняет в базу и пополняет очередь
          найденными ссылками;
    \item \texttt{src/recrawler.py} --- переобкачка: проверяет, не изменились ли уже
          скачанные документы, и изменившиеся ставит в очередь.
\end{itemize}

Ещё три утилиты ничего не качают из сети: \texttt{src/stats.py} печатает сводку по
базе, \texttt{src/validate.py} проверяет, что в ней лежит то, что должно
(раздел~\ref{sec:tests}), а \texttt{src/reparse.py} заново выделяет текст из
сохранённого \texttt{raw\_html}, когда меняются правила разбора.

Общий код разнесён по модулям, у каждого одна ответственность:

\begin{center}
\begin{tabular}{|l|p{11cm}|}
\hline
\textbf{Модуль} & \textbf{Что делает} \\
\hline
\texttt{config.py}   & чтение yaml-конфига \\
\texttt{urls.py}     & нормализация URL \\
\texttt{parsing.py}  & определение кодировки, выделение текста, ссылок и заголовка, отпечаток текста \\
\texttt{rules.py}    & фильтры источника: что обходить, откуда собирать ссылки, что сохранять \\
\texttt{fetcher.py}  & HTTP-клиент: паузы вежливости, \texttt{robots.txt}, условные запросы \\
\texttt{storage.py}  & MongoDB: очередь URL и документы \\
\texttt{shutdown.py} & мягкая остановка по \texttt{SIGINT}/\texttt{SIGTERM} \\
\texttt{stats.py}    & сводка по базе, отдельная утилита \\
\texttt{validate.py} & проверка содержимого базы, отдельная утилита \\
\texttt{reparse.py}  & переразбор текста из сохранённого HTML без перекачки \\
\hline
\end{tabular}
\end{center}

Хранилище и HTTP-клиент передаются в робота и переобкачку снаружи. В рабочем
режиме это настоящие \texttt{Storage} и \texttt{Fetcher}, в автотестах ---
временная база и подменённая сеть (раздел~\ref{sec:tests}).

\subsection*{Хранилище}

В качестве базы данных выбрана MongoDB --- её называет и задание. Документная
модель здесь естественна: документ корпуса --- это запись с набором полей, без
связей между таблицами, а схема легко расширяется новыми полями без миграций.
База поднимается в докере одной командой (\texttt{make db-up}), данные лежат в
именованном томе и переживают пересоздание контейнера. Адрес сервера и имя базы
задаются в секции \texttt{db} конфига.

В базе две коллекции. \texttt{documents} хранит сами документы:

\begin{itemize}
    \item \texttt{url} --- нормализованный адрес, уникальный индекс;
    \item \texttt{raw\_html} --- \enquote{сырой} html-текст документа, как пришёл;
    \item \texttt{text} --- выделенный текст без навигационной обвязки;
    \item \texttt{source} --- название источника;
    \item \texttt{crawled\_at} --- дата обкачки, Unix timestamp;
    \item \texttt{content\_hash} --- отпечаток текста, нужен переобкачке;
    \item \texttt{etag}, \texttt{last\_modified} --- для условных запросов;
    \item \texttt{revision} --- сколько раз документ скачивался;
    \item \texttt{raw\_size}, \texttt{text\_size}, \texttt{text\_bytes} ---
          размеры HTML и текста (в символах и в байтах), чтобы статистика не
          читала тела документов.
\end{itemize}

Задание требует хранить \enquote{сырой} HTML, а лекция --- только необходимый
текст. Требования не противоречат друг другу, если хранить и то, и другое:
\texttt{raw\_html} пригодится, если правила выделения текста изменятся, а
переразбирать корпус по ним, не выкачивая его заново, гораздо дешевле.

Вторая коллекция, \texttt{frontier}, --- это очередь URL: адрес, источник,
глубина, статус (\texttt{pending}, \texttt{done}, \texttt{failed},
\texttt{skipped}), приоритет, число попыток и текст ошибки. Индекс
\texttt{(status, priority, \_id)} делает выбор следующего адреса одним запросом
по индексу. Порядок внутри приоритета --- по \texttt{\_id}: \texttt{ObjectId}
начинается с метки времени, так что это порядок добавления.

Ссылки со страницы добавляются в очередь одной пачкой, \texttt{insert\_many} без
упорядочивания: уже известные адреса отбивает уникальный индекс, остальные
вставляются. Документ сохраняется одной операцией
\texttt{find\_one\_and\_update} с \texttt{upsert}: она же возвращает прежний
отпечаток, по которому видно, изменился ли текст.

\subsection*{Остановка и продолжение}

Очередь лежит в базе, а не в памяти. URL помечается как \texttt{done} только
после того, как документ сохранён, поэтому прерывание в любой момент безопасно:
недообработанный адрес остаётся в статусе \texttt{pending} и при следующем
запуске будет взят снова. Сигналы \texttt{SIGINT} и \texttt{SIGTERM}
перехватываются --- робот дописывает текущий документ и выходит, а не умирает
посреди записи.

На полной обкачке это пригодилось на деле: робот останавливали вручную,
ноутбук засыпал и один раз перезагрузился. Каждый раз следующий запуск
продолжал с того же места.

Уникальность \texttt{url} в обеих коллекциях снимает вопрос дублирования ссылок:
повторное добавление просто игнорируется базой.

\subsection*{Нормализация URL}

Часть ссылок на странице относительные, поэтому при разборе они склеиваются с
адресом текущей страницы. HTML-сущности в атрибуте \texttt{href} раскрываются
(\texttt{\&amp;} --- это \texttt{\&}). Дальше из адреса убирается всё, что не
влияет на содержимое: якорь, служебные параметры (\texttt{utm\_*}, \texttt{gclid}
и подобные), порт по умолчанию, регистр схемы и хоста; параметры запроса
сортируются.

Отдельно пришлось раскрывать процентные последовательности. Ссылки в HTML
Википедии закодированы, и адрес \texttt{/wiki/\%D0\%A4\%D0\%B8\%D0\%BB\%D1\%8C\%D0\%BC}
--- это тот же документ, что и \texttt{/wiki/Фильм}. Без раскрытия в базе
появлялись бы два документа вместо одного. Не раскрываются только \texttt{\%25}
(сам знак процента) и экранированные разделители \texttt{\%2F}, \texttt{\%3F},
\texttt{\%23}: раскрыв их, получим другой адрес.

Повторные слэши в пути сохраняются. Первая версия их схлопывала, и статьи с
названиями вроде \texttt{.hack//SIGN} получали несуществующий адрес и ответ 404.

\subsection*{Определение кодировки}

Порядок ровно такой, как на лекции: сначала заголовок \texttt{Content-Type},
затем \texttt{<meta charset>} в самом HTML, и только потом UTF-8 по умолчанию.
Оба итоговых источника отдают UTF-8, но первоначальный второй источник,
Кино-Театр.Ру, отдавал windows-1251, и без определения кодировки текст с него
попадал бы в базу нечитаемым. Этот случай закреплён в автотестах на сохранённых
страницах сайта.

\subsection*{Вежливость}

Робот соблюдает две вещи. Явные правила --- \texttt{robots.txt}: он скачивается
для каждого хоста и кэшируется, запрещённые адреса помечаются как
\texttt{skipped}. Неявные --- пауза между запросами к одному хосту: общая
задаётся в секции \texttt{logic} конфига (1 секунда), у источника может быть
своя (у film.ru --- 2 секунды, сайт меньше Википедии).

\subsection*{Выбор следующего URL}

Очередь отдаёт сначала адреса, которые ждут переобкачки, из любого источника.
Остальные --- в порядке добавления, то есть обход идёт в ширину.

Из какого источника взять новый адрес, робот решает по готовности хоста:
берётся тот, у которого пауза вежливости истекла раньше, при равенстве --- тот,
к которому дольше не обращались. Это упрощённая версия очереди на каждый хост
из схемы URL frontier.

Первая версия просто чередовала источники, и это оказалось ошибкой. У одного
источника пауза была 5 секунд, у Википедии --- одна, и при строгом чередовании
каждый второй запрос ждал пять секунд: Википедия качалась со скоростью
медленного источника, около 0,4 страницы в секунду. С выбором по готовности
медленный хост получает запрос раз в несколько секунд, а в промежутках робот
качает быстрый.

\subsection*{Сбои сети и блокировки}

Ошибки скачивания делятся на два вида. Настоящие (\texttt{404} и другие ответы
4xx) означают, что страницы нет, --- адрес сразу помечается \texttt{failed}.
Временные --- нет соединения, таймаут, ошибка DNS, ответы 5xx и 429 --- говорят
о том, что страница, скорее всего, есть, но сейчас недоступна. Для них адрес
возвращается в очередь, а хост ставится на паузу: 30 секунд, затем вдвое
дольше, но не больше 10 минут. Пометку \texttt{failed} адрес получает только
после пяти таких сбоев подряд.

Так же обрабатывается капча: если в тексте страницы найден признак заглушки
(\texttt{blocked\_marker} в конфиге источника), страница не сохраняется,
источник встаёт на паузу от 10 минут до 2 часов, а адрес ждёт повторной попытки.

Первая версия любую ошибку считала окончательной, и это стоило дорого: за
несколько часов без интернета робот \enquote{пролистал} очередь и списал как
недоступные 12 тысяч адресов, которые на деле были в порядке (раздел
\enquote{Результаты}).

\subsection*{Выделение текста}

Навигационная обвязка удаляется за счёт того, что текст берётся только из
основного контейнера страницы. Контейнер задаётся в конфиге источника: у
Википедии это \texttt{<div class="mw-content-ltr mw-parser-output">}, у film.ru ---
\texttt{<div class="wrapper\_articles\_text">}. Граница контейнера находится по
балансу вложенных \texttt{<div>}: всё, что после парного \texttt{</div>}, ---
боковые панели, комментарии и подвал сайта. Внутри контейнера выбрасываются
скрипты, стили и таблицы (инфобоксы), затем снимаются теги, а HTML-сущности
раскрываются (\texttt{\&laquo;} --- это \enquote{«}, \texttt{\&amp;} --- \enquote{\&}). Хвост
документа отрезается по маркерам \texttt{cut\_after}: у Википедии это разделы
\enquote{Примечания}, \enquote{Литература} и \enquote{Ссылки}, у film.ru ---
подпись автора и кнопки \enquote{Поделиться}.

\subsection*{Фильтры ссылок}

Фильтров три, и это оказалось важным. Первый, \texttt{allow}/\texttt{deny},
решает, обходить ли ссылку вообще. Второй, \texttt{store\_allow} и
\texttt{store\_deny}, решает, сохранять ли страницу как документ.

Разделять их пришлось потому, что каталоги, списки и страницы категорий нужны
роботу только как источник ссылок. Текста в них хватает, чтобы пройти фильтр по
длине, и в первой версии робота база наполнилась оглавлениями вместо статей.

Третий, \texttt{follow\_only}, --- список страниц, с которых вообще собираются
ссылки. Без него робот уходит из своей тематики: из статьи о фильме ссылки
ведут на произвольные статьи Википедии, и через несколько шагов корпус
перестаёт быть кинематографическим.

\subsection*{Как устроен обход источников}

\textbf{Википедия.} Документы ищутся через категории: робот собирает ссылки
только со страниц категорий, а статьи сохраняет. Сначала стартовыми были
алфавитные категории (\enquote{Фильмы по алфавиту} и подобные), но страница
категории показывает только первые 200 статей. Продолжение лежит по адресам
\texttt{/w/index.php?\ldots\&pagefrom=\ldots}, а весь путь \texttt{/w/}
закрыт в \texttt{robots.txt} Википедии. Поэтому корпус набирается иначе: робот
обходит всё дерево категорий о кино --- по годам, странам, жанрам и их
пересечениям (\enquote{Фильмы Франции 1995 года}). Категорий тысячи, из каждой
берутся первые 200 статей, и в сумме это десятки тысяч документов.

Чтобы обход не ушёл в книги и картины, разрешены только категории, в названии
которых есть \enquote{фильм}, \enquote{сериал}, \enquote{актёр},
\enquote{актриса} или \enquote{кинорежиссёр}; служебные и мета-категории
(\enquote{Википедия:\ldots}, \enquote{Категории\_\ldots}) отсекаются. Отсекаются
и адреса Википедии с любыми параметрами: это ссылки на правку, историю и
\enquote{красные ссылки} на несуществующие страницы (\texttt{?action=edit\&redlink=1}).

\textbf{film.ru.} Статьи собраны в три списка --- рецензии, подборки, интервью,
по 28 статей на странице. Списки листаются через ajax-запрос, который
возвращает JSON, а ссылок на \texttt{?page=N} в HTML нет. Поэтому страницы
списков задаются в конфиге шаблоном \texttt{seed\_pages}:
\texttt{https://www.film.ru/topic/review?page=\{n\}} для $n$ от 0 до последней
страницы. Номер последней страницы найден перебором: следующая отвечает 404.
Робот обходит эти 545 страниц и сохраняет только сами статьи
(\texttt{/articles/\ldots}); короткие карточки фильмов документами не считаются.

\pagebreak
